\subsection{The unit-connector eighth moment for odd terminals with a unit arm and multiple pairs}\label{proof:Q5-C067}

\paragraph{Q67.1 Directions and bounds.}
Use the finite probability space, real skew kernel \(K\), and
operators and messages of Appendix~\ref{proof:Q5-C064}:
\(A=\allowbreak J+K\), \(T=\allowbreak J-K=\allowbreak A^*\), \(g=\allowbreak K\one\), \(p=\allowbreak 1+g\), \(q=\allowbreak 1-g\),
\(f_n=\allowbreak A^n\one\), \(\bar f_n=\allowbreak T^n\one\), \(h_n=\allowbreak Tf_{n-1}\),
\(\delta_j=\allowbreak \|K^j\one\|_2^2\), \(d=\allowbreak \delta_1\),
\(\theta=\allowbreak 13/32\), \(\kappa=\allowbreak 19/32\), and \(\tau=\allowbreak d-\delta_2\).
Fix any finite \(k\ge\allowbreak 2\) and strictly increasing list
\(1<\allowbreak b_1<\allowbreak c_1<\allowbreak b_2<\allowbreak c_2<\allowbreak \cdots<\allowbreak b_k<\allowbreak c_k\).
Let \(E_0,\allowbreak E_r\) be arbitrary finite genuine equal-pair products,
with any lengths and multiplicities; the empty product is one.
Set
\begin{equation}\tag{Q67.1.1}\label{eq:Q67:1.1}
F=E_0qf_{b_1}h_{c_1}\prod_{i=2}^kh_{b_i}f_{c_i},
\qquad H_D=E_rpTF.
\end{equation}
The terminal unit arm points inward. In the first pair the short arm
is outward, while the long arm has its first edge inward and its
remaining edges outward. Every later pair has the short arm first
inward and then outward, and the long arm outward. The unit connector
points from terminal to root; the root unit arm points outward.
Each equal pair consists of an outward arm and an arm first inward
and then outward. These directions depend only on the arm list.
For all such parameters,
\begin{equation}\tag{Q67.1.2}\label{eq:Q67:1.2}
\boxed{\E H_D^8\le e^{-d/2}\le1-2d/5.}
\end{equation}
On \(d\ge\allowbreak 1/24\), the stronger bound is \(e^{-9d}\).
Combining the multiple-pair result with the triple bounds of
Q5-C062, Q5-C063 and Appendix~\ref{proof:Q5-C065} gives, for every
odd distinct residual list of at least three arms and minimum one,
in these same directions,
\begin{equation}\tag{Q67.1.3}\label{eq:Q67:1.3}
\boxed{\E H_D^8\le1-d/601.}
\end{equation}
Three residual arms use the triple bounds; at least five use
\eqref{eq:Q67:1.2}. Residual zero and one follow the reduction to
fewer branch vertices.

The finite convolution, the all-index mean bound
\(\mu_n\le\allowbreak 1-83d/32\) for \(n\ge\allowbreak 4\), and the finite-mask norm are
proved in Appendix~\ref{proof:Q5-C090}.
Its long-pair estimate \eqref{eq:L90:pair} supplies in particular
\(h_af_c\le\allowbreak  e^{R_{ac}-3d}\) for \(4\le\allowbreak  a<\allowbreak c\), where
\(R_{ac}=\allowbreak \sum_{j=\allowbreak a}^{c-1}K^jg\).
We use the finite identity
\(S_n=\allowbreak f_n+\sum_{j=\allowbreak 3}^n\beta_jf_{n-j}\), with
\(S_n=\allowbreak \sum_{j=\allowbreak 0}^nK^j\one\),
\(\beta_1=\allowbreak \beta_2=\allowbreak 0\),
\(\beta_n=\allowbreak \sum_{j=\allowbreak 1}^{\lfloor(n-1)/2\rfloor}(-1)^{j-1}\delta_j\),
and \(\tau\le\allowbreak \beta_n\le\allowbreak  d\) for \(n\ge\allowbreak 3\).

\paragraph{Q67.3 Unit propagation retaining the negative mean.}
Suppose \(0<\allowbreak d\le\allowbreak 1/24\) and a nonnegative message \(F\) satisfies
\begin{equation}\tag{Q67.3.1}\label{eq:Q67:3.1}
F\le e^{R-\Gamma},\quad |R|\le5/4,\quad
\E R^2\le A_0d,\quad\E(gR)\le0,\qquad A_0>0.
\end{equation}
Put \(N=\allowbreak -KR\). Row and spectral bounds, followed by Cauchy in the
same inner product, give
\begin{equation}\tag{Q67.3.2}\label{eq:Q67:3.2}
|N|\le\sqrt{A_0d},\quad \E N^2\le\theta A_0d,\quad
\E N=\E(gR)\le0,\quad|\E N|\le\sqrt{A_0}d.
\end{equation}
The negative-mean lemma \eqref{eq:Q65:2.4}, including its two-point
comparison, Hermite interpolation and degenerate cases, gives
\begin{equation}\tag{Q67.3.3}\label{eq:Q67:3.3}
\E e^{8N}\le\exp\{8\E N+\lambda(\sqrt d)A_0d\},\qquad
\lambda(x)=\frac{188}{25}\left(2-\frac{\kappa}{1+x}\right)^2.
\end{equation}
Since \(\lambda\) increases and \(\sqrt d\le\allowbreak 5/24\),
\begin{equation}\tag{Q67.3.4}\label{eq:Q67:3.4}
\lambda(\sqrt d)\le\lambda(5/24)=57575/3364<69/4,
\end{equation}
where \(230300<\allowbreak 232116\).
For every real \(u\le\allowbreak 5/4\), one has \(e^u\le\allowbreak 1+u+(4/5)u^2\).
Indeed,
\((e^u-1-u)/u^2=\allowbreak \int_0^1(1-t)e^{tu}\,\allowbreak dt\) increases with \(u\),
so it suffices that \(e^{5/4}<\allowbreak 7/2\).
The exponential series through degree four sums to \(21249/6144\);
starting with degree five, adjacent tail terms have ratio at most
\(5/24\). Thus the total is at most
\fitdisplay{21249/6144+3750/116736=407481/116736<7/2.}{}
Positivity of \(T\) and \(|K(R^2)|\le\allowbreak \E R^2\) give
\begin{equation}\tag{Q67.3.5}\label{eq:Q67:3.5}
Te^R\le q+\E R-KR+(8/5)\E R^2=q+C+N,
\qquad C=\E R+(8/5)\E R^2.
\end{equation}
The right side is nonnegative. Since \(p+q=\allowbreak 2\), AM--GM gives
\(p(q+C+N)\le\allowbreak [1+(C+N)/2]^2\) with nonnegative base.
Use \((1+v)^{16}\le\allowbreak  e^{16v}\) for \(v\ge\allowbreak -1\), obtaining
\begin{equation}\tag{Q67.3.6}\label{eq:Q67:3.6}
\boxed{\E[(pTF)^8]\le
\exp\{-8\Gamma+8\E R+8\E(gR)+(601/20)A_0d\}.}
\end{equation}
Here \(601/20=\allowbreak 64/5+69/4\). A zero base satisfies the comparison
directly. Both signed means remain in this bound.

\paragraph{Q67.4 The odd-table envelope and high domain.}
First take \(E_0=\allowbreak E_r=\allowbreak 1\). Write \(b=\allowbreak b_1\), \(c=\allowbreak c_1\),
\(B=\allowbreak \beta_{c-1}\), and define
\begin{equation}\tag{Q67.4.1}\label{eq:Q67:4.1}
\begin{gathered}
W=\sum_{j=b}^{c-1}K^jg,\qquad
S=\sum_{i=2}^k\sum_{j=b_i}^{c_i-1}K^jg,\\
\Delta=\beta_b+\beta_c+B,\quad\alpha=1-B,\quad
R=-\alpha g-W+S,\quad\Gamma=\Delta+3(k-1)d.
\end{gathered}
\end{equation}
Finite convolution and \(h_c=\allowbreak Tf_{c-1}\) give
\(f_b+h_c\le\allowbreak 2-\Delta+Bg-W\).
Apply AM--GM, \(q\le\allowbreak  e^{-g}\), and the long-pair bound to obtain
\begin{equation}\tag{Q67.4.2}\label{eq:Q67:4.2}
qf_bh_c\prod_{i=2}^kh_{b_i}f_{c_i}\le e^{R-\Gamma}.
\end{equation}
The later exponent intervals are disjoint, so
\(S=\allowbreak K^{c+1}Q(K)g\), where \(Q\) has coefficients zero or one and
may have initial zero coefficients. The finite-mask bound is
\begin{equation}\tag{Q67.4.3}\label{eq:Q67:4.3}
\|Q(K)\|^2\le\frac{1+\theta}{(1-\theta^2)^2}
=32768/16245<101/50<(10/7)^2.
\end{equation}
For \(c\ge\allowbreak 3\), this implies
\(|S|\le\allowbreak (10/7)\theta^{3/2}\sqrt d<\allowbreak (3/8)\sqrt d\);
the squared comparison is \(14060800<\allowbreak 14450688\).
If \(d\ge\allowbreak 1/24\), \(\sqrt d\le\allowbreak \sqrt{24}d<\allowbreak 5d\), so
\begin{equation}\tag{Q67.4.4}\label{eq:Q67:4.4}
\prod_{i=2}^kh_{b_i}f_{c_i}
\le e^{S-3(k-1)d}\le e^{-9d/8}.
\end{equation}
Positivity and the same-direction unit eighth moments of Q5-C062
for \((b,\allowbreak c)=\allowbreak (2,\allowbreak 3)\), Q5-C063 for \(b=\allowbreak 2,\allowbreak c\ge\allowbreak 4\), and
Appendix~\ref{proof:Q5-C065} for \(b\ge\allowbreak 3\), give
\begin{equation}\tag{Q67.4.5}\label{eq:Q67:4.5}
\E[(pTF)^8]\le e^{-9d}\E[(pT(qf_bh_c))^8]\le e^{-9d}.
\end{equation}

\paragraph{Q67.5 Signs, pointwise bounds and finite frequencies on the low domain.}
Now let \(0<\allowbreak d\le\allowbreak 1/24\), so \(23/24\le\allowbreak \alpha\le\allowbreak 1\).
Every finite \(G_m\) satisfies \(\|G_m(K)\|\le\allowbreak \sqrt{1+\theta}<\allowbreak 6/5\).
Apply the original row bound to the outermost \(K\) in \(W\) and \(S\):
\begin{equation}\tag{Q67.5.1}\label{eq:Q67:5.1}
|R|\le1+\left((6/5)\sqrt\theta+(10/7)\theta^{3/2}\right)\sqrt d
<1+(6/5)\sqrt d\le5/4.
\end{equation}
The middle coefficient is bounded using \(\sqrt\theta<\allowbreak 2/3\),
which gives \(997/840<\allowbreak 6/5\).
For the single interval \(W\), alternating sums give
\(-\E(gW)\le\allowbreak \delta_2\). The positive mixed-mean terms of \(S\)
begin at \(K^4g\), and their sum is at most
\(\delta_3/(1-\theta^2)<\allowbreak \delta_2/2\). Therefore
\begin{equation}\tag{Q67.5.2}\label{eq:Q67:5.2}
\E(gR)\le-\alpha d+(3/2)\delta_2\le-67d/192<0.
\end{equation}
We next prove, for all finite intervals and masks,
\begin{equation}\tag{Q67.5.3}\label{eq:Q67:5.3}
\E R^2\le A_0d,\qquad
A_0=\begin{cases}
1,&b=2,\\33/20,&b=3,\\19/12,&b=4,\\5/4,&b\ge5.
\end{cases}
\end{equation}
On a skew-operator rotation block, put \(z=\allowbreak i\sqrt x\),
\(0\le\allowbreak  x\le\allowbreak \theta\), and write \(W=\allowbreak W_0(K)g\), \(S=\allowbreak U(K)g\).
The polynomial representing \(-R\) is
\(P(z)=\allowbreak \alpha+W_0(z)-U(z)\), where
\(W_0=\allowbreak \sum_{j=\allowbreak b}^{c-1}z^j\), and the zero-one coefficients of \(U\)
begin at degree at least \(c+1\). Finite sums in the four residue
classes give the coordinate bounds below. We use
\begin{equation}\tag{Q67.5.4}\label{eq:Q67:5.4}
\begin{gathered}
\frac{x^2}{1-x^2}\le169/855<1/5,\quad
\frac{x}{1-x^2}\le416/855<1/2,\quad
\frac{x^3}{1-x^2}<1/12,\\
x^2<1/6,\quad x^3<1/14,\quad x^4<1/32,\quad
x^5<1/90,\quad\frac{x^4}{1-x^2}<1/30.
\end{gathered}
\end{equation}
The third fraction is at most \(2197/27360<\allowbreak 1/12\).
For \(x^3\) and \(x^5\), the integer comparisons are
\(30758<\allowbreak 32768\) and \(33416370<\allowbreak 33554432\); the remaining bounds
follow by cross multiplication.

\paragraph{Q67.5.1 The class \(b=\allowbreak 2,\allowbreak c\ge\allowbreak 3\).}
At \(c=\allowbreak 3\), \(B=\allowbreak 0\), \(P=\allowbreak 1+z^2-U\), with \(U\) starting at degree
four. The real part is positive and at most \(1-4x/5\), and the
imaginary magnitude is at most \(x^{5/2}/(1-x^2)\). Hence
\fitdisplay{
\frac{|P|^2-1}{x}\le-8/5+16x/25+\frac{x^4}{(1-x^2)^2}
\le-8/5+8/25+1/25<0.
}{}
At \(c=\allowbreak 4\), \(W_0=\allowbreak z^2+z^3\), with \(U\) starting at degree five.
The real part is positive and at most \(1-4x/5\); the imaginary
part is negative, with magnitude at most \((3/2)x^{3/2}\). Thus
\fitdisplay{
\frac{|P|^2-1}{x}\le-8/5+16x/25+9x^2/4
\le-8/5+8/25+3/8<0.
}{}
For \(c\ge\allowbreak 5\), the real part of \(W_0\) lies in
\([-x,\allowbreak -x+x^2]\), and its imaginary part in
\([-x^{3/2},\allowbreak -x^{3/2}(1-x)]\).
The mask \(U\) starts at least at degree six, with positive imaginary
terms starting at nine and negative imaginary terms at seven.
The real part of \(P\) is positive and
\fitdisplay{\operatorname{Re}P\le1-x+x^2+\frac{x^3}{1-x^2}
\le1-3x/8.}{}
Here \(1-x-x^2/(1-x^2)\ge\allowbreak 1-\theta-1/5=\allowbreak 63/160>\allowbreak 3/8\).
Its imaginary part remains nonpositive because
\(1-x>\allowbreak x^2/(1-x^2)\), and
\fitdisplay{|\operatorname{Im}P|\le
x^{3/2}\left(1+\frac{x^3}{1-x^2}\right)
\le(13/12)x^{3/2}.}{}
Consequently
\fitdisplay{
\frac{|P|^2-1}{x}\le-3/4+9x/64+169x^2/144
\le-3/4+9/128+169/864=-1673/3456<0.
}{}
The real-part positivity used above follows from
\(1-x-x^2/(1-x^2)>\allowbreak 0\) or
\(23/24-x-x^4/(1-x^2)>\allowbreak 0\), respectively. At zero the bound is direct.
This proves \(A_0=\allowbreak 1\) for every finite \(c\).

\paragraph{Q67.5.2 The class \(b=\allowbreak 3,\allowbreak c\ge\allowbreak 4\).}
At \(c=\allowbreak 4\), \(W_0=\allowbreak z^3\), and \(U\) starts at degree five.
The real part is positive and at most \(13/12\), and the imaginary
magnitude is at most \((3/2)x^{3/2}\). Thus
\(|P|^2\le\allowbreak 169/144+9/56=\allowbreak 1345/1008<\allowbreak 33/20\).
For \(c\ge\allowbreak 5\), \(\operatorname{Re}W_0\in[0,\allowbreak x^2]\),
\(\operatorname{Im}W_0\in[-x^{3/2},\allowbreak -x^{3/2}(1-x)]\).
The mask starts at least at degree six, with negative real terms
starting at six and positive imaginary terms at nine.
The real part of \(P\) is positive and at most
\(1+x^2+x^3/(1-x^2)\le\allowbreak 5/4\).
As above its imaginary part is nonpositive, of magnitude at most
\((13/12)x^{3/2}\). Therefore
\fitdisplay{|P|^2\le25/16+169/2016=3319/2016<33/20,}{}
where \(66380<\allowbreak 66528\).

\paragraph{Q67.5.3 The class \(b=\allowbreak 4,\allowbreak c\ge\allowbreak 5\).}
Here \(\operatorname{Re}W_0\in[x^2(1-x),\allowbreak x^2]\) and
\(\operatorname{Im}W_0\in[0,\allowbreak x^{5/2}]\); the latter is zero at
\(c=\allowbreak 5\). The mask begins at least at degree six.
The real part is at least \(23/24-x^4/(1-x^2)>\allowbreak 0\), and at most
\fitdisplay{1+x^2+\frac{x^3}{1-x^2}\le319/256.}{}
This uses \(\theta^2=\allowbreak 169/1024\) and
\(2197/27360<\allowbreak 83/1024\), with \(2249728<\allowbreak 2270880\).
The imaginary magnitude is at most
\(x^{5/2}[1+x/(1-x^2)]\le\allowbreak (3/2)x^{5/2}\). Thus
\begin{equation}\tag{Q67.5.5}\label{eq:Q67:5.5}
|P|^2\le101761/65536+1/40<19/12,
\end{equation}
where the final comparison is \(12211320<\allowbreak 12255232\).

\paragraph{Q67.5.4 The class \(b\ge\allowbreak 5,\allowbreak c>\allowbreak b\).}
The contiguous \(W_0\) has real part in \([-x^3,\allowbreak x^4]\) and imaginary
magnitude at most \(x^{5/2}\).
The mask begins at least at degree seven: its negative real terms
start at ten, its positive real terms at eight, and its imaginary
magnitude is at most \(x^{7/2}/(1-x^2)\).
Thus the real part of \(P\) is positive and
\fitdisplay{\operatorname{Re}P\le1+x^4+\frac{x^5}{1-x^2}
<1+1/32+1/70<21/20.}{}
The imaginary magnitude is at most \((3/2)x^{5/2}\).
For positivity use \(23/24-x^3-x^4/(1-x^2)>\allowbreak 0\).
It follows that \(|P|^2\le\allowbreak 441/400+1/40=\allowbreak 451/400<\allowbreak 5/4\).
These four classes prove \eqref{eq:Q67:5.3} for arbitrary finite
degree, pair count and gaps.

\paragraph{Q67.6 The common linear functional and continuous low-domain bounds.}
Write \(a_j=\allowbreak \E K^j\one\), so \(a_{2j}=\allowbreak (-1)^j\delta_j\) and every
odd \(a_j\) is zero. Skew adjunction gives the telescoping identity
\begin{equation}\tag{Q67.6.1}\label{eq:Q67:6.1}
\E W+\E(gW)=a_{b+1}-a_{c+1}.
\end{equation}
Put \(V=\allowbreak a_{c+1}-a_{b+1}\). Then
\begin{equation}\tag{Q67.6.2}\label{eq:Q67:6.2}
\E R+\E(gR)=-\alpha d+V+\E S+\E(gS).
\end{equation}
For \((b,\allowbreak c)=\allowbreak (2,\allowbreak 3)\), \(S\) starts at degree four. The positive
\(\E S\) terms begin at degree seven with \(\delta_4\), while the
positive \(\E(gS)\) terms begin at degree four with \(\delta_3\).
Therefore
\begin{equation}\tag{Q67.6.3}\label{eq:Q67:6.3}
\E S+\E(gS)\le\frac{\delta_3+\delta_4}{1-\theta^2}
\le\frac{\theta^2}{\kappa}d=169d/608<7d/25.
\end{equation}
For the other branches \(c\ge\allowbreak 4\), the mask starts at least at degree
five. The positive \(\E S\) terms start at seven with \(\delta_4\),
and the positive \(\E(gS)\) terms at eight with \(\delta_5\). Thus
\begin{equation}\tag{Q67.6.4}\label{eq:Q67:6.4}
\E S+\E(gS)\le\frac{\delta_4+\delta_5}{1-\theta^2}
\le\frac{\theta^3}{\kappa}d=2197d/19456<23d/200.
\end{equation}
The integer comparisons are \(4225<\allowbreak 4256\) and \(439400<\allowbreak 447488\).
Discarding the negative terms of the finite mask gives these upper
bounds for the common moments.
Equations~\eqref{eq:Q67:3.6}, \eqref{eq:Q67:4.1} and
\eqref{eq:Q67:6.2} give
\begin{equation}\tag{Q67.6.5}\label{eq:Q67:6.5}
\begin{aligned}
\log\E[(pTF)^8]\le{}&-8(\Delta-V)-8d+8Bd-24(k-1)d\\
 &+8(\E S+\E(gS))+(601/20)A_0d.
\end{aligned}
\end{equation}
At zero expectation, take the logarithm to be \(-\infty\);
the tree messages on positive-mass atoms are positive.
All following bounds hold on \(0<\allowbreak d\le\allowbreak 1/24\), using \(k\ge\allowbreak 2\).

For \((b,\allowbreak c)=\allowbreak (2,\allowbreak 3)\), \(\Delta=\allowbreak d\), \(B=\allowbreak 0\), \(V=\allowbreak \delta_2\).
Use \eqref{eq:Q67:6.3} and \(A_0=\allowbreak 1\). Dividing the upper exponent
by \(d\) gives
\begin{equation}\tag{Q67.6.6}\label{eq:Q67:6.6}
-40+8\theta+56/25+601/20=-223/50<-1/2.
\end{equation}
For \(b=\allowbreak 2,\allowbreak c\ge\allowbreak 4\), \(\Delta\ge\allowbreak 2\tau\) and
\(V=\allowbreak a_{c+1}\le\allowbreak \delta_4\), since the first positive even moment for
\(c+1\ge\allowbreak 5\) is at \(c+1=\allowbreak 8\). Thus \(\Delta-V\ge\allowbreak 2\tau-\delta_4\).
With \eqref{eq:Q67:6.4}, \(8B\le\allowbreak 1/3\), and \(A_0=\allowbreak 1\), the coefficient is at most
\begin{equation}\tag{Q67.6.7}\label{eq:Q67:6.7}
-48+16\theta+8\theta^3+1/3+23/25+601/20<-9<-1/2.
\end{equation}
Use \(16\theta=\allowbreak 13/2\) and \(8\theta^3<\allowbreak 4/7\) for the last bound.
For \(b=\allowbreak 3\), \(\Delta\ge\allowbreak  d+2\tau\) and
\(V=\allowbreak a_{c+1}-\delta_2\le\allowbreak \delta_4-\delta_2\).
Hence \(\Delta-V\ge\allowbreak 3d-\delta_2-\delta_4\), and \(A_0=\allowbreak 33/20\) gives
\begin{equation}\tag{Q67.6.8}\label{eq:Q67:6.8}
-56+8\theta+8\theta^3+1/3+23/25+19833/400
<-11279/8400<-1<-1/2.
\end{equation}
Substituting \(8\theta=\allowbreak 13/4\), \(8\theta^3<\allowbreak 4/7\) gives the stated
\(-11279/8400\).
For \(b=\allowbreak 4\), \(V=\allowbreak a_{c+1}\). At \(c=\allowbreak 5,\allowbreak 6\), \(V\le\allowbreak 0\) and
\(\Delta\ge\allowbreak 3d-2\delta_2\).
For \(c\ge\allowbreak 7\), \(\beta_c\ge\allowbreak \tau+\kappa\delta_3\),
\(\beta_{c-1}\ge\allowbreak \tau\), and \(V\le\allowbreak \delta_4\le\allowbreak \theta\delta_3\).
Since \(\kappa>\allowbreak \theta\), this gives throughout
\begin{equation}\tag{Q67.6.9}\label{eq:Q67:6.9}
\Delta-V\ge3d-2\delta_2.
\end{equation}
The extra term in \(\beta_c\) is the alternating tail
\(\delta_3-\delta_4+\cdots\ge\allowbreak \delta_3-\delta_4\), with precisely
\(+\delta_3\) at \(c=\allowbreak 7\). This preserves the common-moment relation.
With \(A_0=\allowbreak 19/12\), the coefficient is at most
\begin{equation}\tag{Q67.6.10}\label{eq:Q67:6.10}
-56+16\theta+1/3+23/25+11419/240=-267/400<-1/2.
\end{equation}
Its margin below \(-1/2\) is \(67/400\).
For \(b\ge\allowbreak 5\), \(\Delta\ge\allowbreak 3\tau\), \(V\le\allowbreak \delta_3+\delta_4\):
the first positive contribution of \(-a_{b+1}\) is \(\delta_3\),
and that of \(a_{c+1}\) is \(\delta_4\).
Using \(A_0=\allowbreak 5/4\), the coefficient is at most
\begin{equation}\tag{Q67.6.11}\label{eq:Q67:6.11}
-56+24\theta+8\theta^2+8\theta^3+1/3+23/25+601/16<-5<-1/2.
\end{equation}
For this bound use \(\theta^2<\allowbreak 1/6\), \(\theta^3<\allowbreak 1/14\).
Equations~\eqref{eq:Q67:6.6}--\eqref{eq:Q67:6.11} prove
\(\E[(pTF)^8]\le\allowbreak  e^{-d/2}\) on the low domain.

\paragraph{Q67.7 Completion and the full odd unit-arm list.}
The high-domain bound \eqref{eq:Q67:4.5} and the low-domain proof
cover \(d>\allowbreak 0\) with the same directions.
At \(d=\allowbreak 0\), \(g=\allowbreak 0\) and \(K\) may be nonzero; leaf integration
makes all path, spider and tree messages one.
Since \(0\le\allowbreak  E_0,\allowbreak E_r\le\allowbreak 1\) and \(T\) is positive,
\fitdisplay{
0\le E_rpT\left(E_0qf_bh_c\prod_{i=2}^kh_{b_i}f_{c_i}\right)
\le pT\left(qf_bh_c\prod_{i=2}^kh_{b_i}f_{c_i}\right).
}{}
Thus all genuine equal-pair decorations satisfy the bound. Moreover,
\begin{equation}\tag{Q67.7.1}\label{eq:Q67:7.1}
e^{-d/2}\le(1+d/2)^{-1}\le1-3d/7\le1-2d/5
\qquad(0\le d\le1/3).
\end{equation}
Apply the theorem to \(-K\), exchanging \(A/T\), \(p/q\) and every
edge together, to obtain the reversed \(qAF^-\) ordinary eighth moment.
The same multiple-pair direction is that of Q4-C138, whose
all-positive-length fourth moment is proved in Section~\ref{loc:S35}:
\fitdisplay{\E[p(T^LF)^4]\le1\qquad(L\ge1).}{}
The triple directions have the fourth moments of Q4-C125, Q4-C129
and Q4-C134. For their unit eighth moments, Q5-C062 gives
\(1-d/601\), and Q5-C063 and Q5-C065 give \(1-d/500\).
Together with \(1-2d/5\) for every list of at least five arms,
these yield \eqref{eq:Q67:1.3}.
